\documentclass[10pt]{article}

% ---------- Document Identity (update these only) ----------
\newcommand{\DocStage}{}
\newcommand{\DocTitle}{Your Road to Tier One SOF}
\newcommand{\ProgramName}{182-Week Civilian-to-Selection Readiness Pipeline}
\newcommand{\ProgramSubtitle}{}

% ---------- Page ----------
\usepackage[legalpaper,left=0.5in,right=0.5in,top=0.6in,bottom=0.45in]{geometry}

% ---------- Typography ----------
\usepackage[T1]{fontenc}
\usepackage[utf8]{inputenc}
\usepackage{libertinus}
\usepackage{microtype}
\usepackage{parskip}
\usepackage{ragged2e}
\usepackage{textcomp}
\AtBeginDocument{\color{black}}
\AtBeginDocument{\pagecolor{white}}
\AtBeginDocument{\justifying}

% ---------- Landscape pages ----------
\usepackage{pdflscape}

% ---------- Color ----------
\usepackage[table]{xcolor}
\definecolor{StageBlue}{HTML}{1F4E79}
\definecolor{StageBlueDark}{HTML}{163A5A}
\definecolor{LightGray}{HTML}{F2F4F7}
\definecolor{MidGray}{HTML}{D9DEE5}
\definecolor{RuleGray}{HTML}{C7CED8}
\definecolor{HeaderInk}{HTML}{000000}
\definecolor{Ink}{HTML}{000000}

% ---------- Utility ----------
\usepackage{enumitem}
\sloppy
\setlength{\emergencystretch}{2em}

% ---------- Boxes / UI ----------
\usepackage[most]{tcolorbox}

\tcbset{
  charterTitle/.style={
    enhanced,
    colback=StageBlue!4,
    colframe=StageBlue,
    boxrule=1.25pt,
    arc=3pt,
    left=16pt,right=16pt,top=14pt,bottom=14pt,
    borderline north={3.2pt}{0pt}{StageBlue},
    borderline west={4pt}{0pt}{StageBlue},
    borderline south={1.25pt}{0pt}{StageBlue},
    drop shadow={black!25!white},
  },
  charterRule/.style={
    enhanced,
    colback=LightGray,
    colframe=RuleGray,
    boxrule=0.85pt,
    arc=3pt,
    left=12pt,right=12pt,top=10pt,bottom=10pt,
    borderline west={3pt}{0pt}{StageBlue},
  },
  charterBadge/.style={
    enhanced,
    colback=StageBlue,
    coltext=white,
    colframe=StageBlueDark,
    boxrule=0.6pt,
    arc=2pt,
    left=6pt,right=6pt,top=3pt,bottom=3pt,
  }
}
\newtcbox{\Badge}{nobeforeafter,charterBadge}

% ---------- Lists ----------
\setlist[itemize]{leftmargin=1.5em, itemsep=0.25em, topsep=0.25em}
\setlist[enumerate]{leftmargin=1.8em, itemsep=0.25em, topsep=0.25em}

% ---------- TOC ----------
\usepackage{tocloft}
\setcounter{tocdepth}{3}
\setlength{\cftbeforesecskip}{3pt}
\setlength{\cftbeforesubsecskip}{2pt}
\setlength{\cftbeforesubsubsecskip}{0pt}
\renewcommand{\cftsecfont}{\bfseries}
\renewcommand{\cftsecpagefont}{\bfseries}
\renewcommand{\cftsubsecfont}{\normalfont}
\renewcommand{\cftsubsecpagefont}{\normalfont}
\renewcommand{\cftsubsubsecfont}{\normalfont}
\renewcommand{\cftsubsubsecpagefont}{\normalfont}
\setlength{\cftsecindent}{0pt}
\setlength{\cftsubsecindent}{1.25em}
\setlength{\cftsubsubsecindent}{2.8em}
\setlength{\cftsecnumwidth}{2.1em}
\setlength{\cftsubsecnumwidth}{2.8em}
\setlength{\cftsubsubsecnumwidth}{3.6em}

% Remove the built-in TOC heading so only the custom heading is shown
\makeatletter
\renewcommand{\tableofcontents}{\@starttoc{toc}}
\makeatother

% ---------- Header/Footer: page number only ----------
\usepackage{fancyhdr}
\pagestyle{fancy}
\fancyhf{}
\renewcommand{\headrulewidth}{0pt}
\renewcommand{\footrulewidth}{0pt}
\fancyfoot[C]{\thepage}

% ---------- Section formatting ----------
\usepackage{titlesec}

\titleformat{\section}
  {\Large\bfseries\color{Ink}}
  {\thesection}{0.6em}{}
  [\vspace{0.25em}\textcolor{StageBlue}{\rule{\linewidth}{0.8pt}}]

\titleformat{\subsection}
  {\large\bfseries\color{Ink}}
  {\thesubsection}{0.6em}{}

\titleformat{\subsubsection}
  {\normalsize\bfseries\color{Ink}}
  {\thesubsubsection}{0.6em}{}

\titlespacing*{\section}{0pt}{0.95em}{0.35em}
\titlespacing*{\subsection}{0pt}{0.65em}{0.25em}
\titlespacing*{\subsubsection}{0pt}{0.55em}{0.2em}

% ---------- Table engine ----------
\usepackage{tabularray}
\UseTblrLibrary{booktabs}

% ---------- tabularray: GLOBAL table treatment ----------
\SetTblrInner{
  cells = {fg=black, bg=white, valign=t},
}

% ---------- Header helpers ----------
\newcommand{\Hdr}[1]{\shortstack[c]{\strut #1\strut}}

% ---------- Lines ----------
\newcommand{\TitleLine}{\textcolor{StageBlue}{\rule{0.98\linewidth}{0.95pt}}}
\newcommand{\SubTitleLine}{\textcolor{RuleGray}{\rule{0.98\linewidth}{0.65pt}}}

% ---------- Continued on next page ----------
\newcommand{\ContinuedOnNextPageInline}{%
  \par
  \vfill
  \noindent\textcolor{RuleGray}{\rule{\linewidth}{1pt}}\vspace{-2pt}
  \noindent\hfill\textit{Continued on next page}\par\vspace{-10pt}
  \noindent\textcolor{RuleGray}{\rule{\linewidth}{1pt}}
  \clearpage
}

% ---------- PDF hyperlinks ----------
\usepackage[hidelinks]{hyperref}
\hypersetup{
  colorlinks=false,
  pdfborder={0 0 0},
  linktoc=all,
  pdftitle={\DocTitle},
  pdfauthor={\ProgramName}
}

% =========================================================
\begin{document}

\section*{7.07 — Load-Bearing, Rucking, and Carry Systems}
\phantomsection
\addcontentsline{toc}{section}{7.07 — Load-Bearing, Rucking, and Carry Systems}

\subsection*{7.07.1 Purpose \& Scope}
\phantomsection
\addcontentsline{toc}{subsection}{7.07.1 Purpose \& Scope}

7.07 defines the load-bearing system stack required to carry external load safely, repeatably, and audibly across phases. It covers ruck pack systems, load modules/media, and carry implements, plus fit/adjustment and safety doctrine needed to prevent predictable injury, foot breakdown, and invalid evidence.

This category does not cover, except by cross-reference:

\begin{itemize}
\item footwear and foot care: 7.11
\item apparel systems: 7.14 and 7.15
\item nutrition/hydration containers and prep systems: 7.13
\item protective/visibility gear: 7.19
\item testing/timing/measurement tooling and load verification method posture: 7.18
\end{itemize}

\subsection*{7.07.2 Governance and Safety Boundaries}
\phantomsection
\addcontentsline{toc}{subsection}{7.07.2 Governance and Safety Boundaries}

\begin{itemize}
\item Start-from-zero posture: early phases remain minimal. Load carriage is phase-gated and does not become required until Stage 6 indicates Tier 1 readiness.
\item No novelty near gates: ruck interface variables (pack model/class, frame type, belt fit approach, load module type, and retention method) are treated as controlled variables. Do not introduce first-time changes inside Deload+Test windows.
\item Stop posture always honored:
  \begin{itemize}
  \item stop if pain alters mechanics (limp, compensation, unstable joint position),
  \item stop for red-flag symptoms per governing posture,
  \item stop and reslot if environment/visibility makes safe execution uncertain.
  \end{itemize}
\item Traffic and visibility rules are mandatory:
  \begin{itemize}
  \item daylight bias where possible,
  \item reflective/visibility posture per 7.19,
  \item no roadside rucking when safe separation cannot be maintained.
  \end{itemize}
\item Heat/ice/surface safety is mandatory:
  \begin{itemize}
  \item substitute indoors or reslot when unsafe footing, lightning/storm risk, unsafe heat, or poor air quality exists.
  \end{itemize}
\item High orthopedic cost acknowledgement:
  \begin{itemize}
  \item load carriage is a high-cost exposure and must be introduced only when Tier 1 capability exists and foot interface is stable (7.11).
  \end{itemize}
\end{itemize}

\subsection*{7.07.3 Tier Doctrine Applied to Ruck and Carry Systems}
\phantomsection
\addcontentsline{toc}{subsection}{7.07.3 Tier Doctrine Applied to Ruck and Carry Systems}

\begin{itemize}
\item Tier 0: no load carriage capability. No substitutions. Education only. Household backpacks are not treated as capability.
\item Tier 1: minimum viable ruck capability for safe progression and basic exposure: fit-adjustable pack with supportive hip belt, secure load module, and anti-shift retention.
\item Tier 2: expanded comfort/stability/modularity that reduces friction and injury risk: better load transfer, more stable retention, modular load cells for staged changes, and optional carry implements.
\item Tier 3: advanced/overbuilt redundancy and specialized systems: backup pack/frame options, expanded modularity, and carry tool ecosystem to reduce variance and failure probability in high-consequence phases.
\end{itemize}

\subsection*{7.07.4 Selection Criteria \& Fit Checks}
\phantomsection
\addcontentsline{toc}{subsection}{7.07.4 Selection Criteria \& Fit Checks}

Fit and safety checklist (operational):

\begin{itemize}
\item Torso length adjustment: pack rides high; shoulder straps support without digging; no excessive shoulder-only load.
\item Hip belt positioning: belt rides on hips; load transfer to hips is achievable; belt does not slip during walking.
\item Strap stability: sternum strap stabilizes without restricting breathing; shoulder straps do not loosen under movement.
\item Load stability: load sits high/close; minimal sway; no audible shifting; retention and compression keep load fixed.
\item Load rating and integrity:
  \begin{itemize}
  \item stitching intact; webbing not frayed; buckles not cracked; attachment points stable.
  \item any structural doubt is a stop condition for load carriage escalation.
  \end{itemize}
\item Chafe/heat management check:
  \begin{itemize}
  \item identify contact points (shoulders, hips, low back) and verify no hotspot escalation.
  \item cross-reference 7.11 for foot hotspots and blister prevention posture.
  \end{itemize}
\item Cross-references:
  \begin{itemize}
  \item hydration routing and containers are 7.13.
  \item reflective/visibility attachment points are 7.19.
  \item rain/cold management is 7.14; general apparel is 7.15.
  \end{itemize}
\end{itemize}

\subsection*{7.07.5 Setup, Safety, and Anchoring}
\phantomsection
\addcontentsline{toc}{subsection}{7.07.5 Setup, Safety, and Anchoring}

Setup doctrine (no programming):

\begin{itemize}
\item Packing doctrine:
  \begin{itemize}
  \item load high and close to the back to reduce sway and lumbar torque.
  \item use compression straps to eliminate movement; avoid low-hanging loads.
  \item label dry load and keep placement consistent for comparability.
  \end{itemize}
\item Route safety:
  \begin{itemize}
  \item prefer routes with low traffic, wide shoulders/paths, and predictable surfaces.
  \item daylight bias where possible; if low light is unavoidable, visibility posture (7.19) is mandatory.
  \item avoid headphones in traffic environments where situational awareness is required.
  \end{itemize}
\item Surface/traction checks:
  \begin{itemize}
  \item do not ruck on ice glaze or unsafe traction surfaces.
  \item reslot to indoor substitutes when outdoor conditions are unsafe.
  \end{itemize}
\item Environmental limits:
  \begin{itemize}
  \item heat: conservative exposure; hydration system posture (7.13); stop for heat illness signs.
  \item cold/wet: apparel posture (7.14) and drying/sustainment posture (7.10/7.20) prevent skin breakdown.
  \end{itemize}
\item Home safety:
  \begin{itemize}
  \item stairs and slick floors are high fall risk with load; avoid if not controlled.
  \item eliminate trip hazards; confirm doorway widths and turning clearance before loaded indoor carries.
  \end{itemize}
\end{itemize}

\subsection*{7.07.6 Maintenance, Replacement, and Storage}
\phantomsection
\addcontentsline{toc}{subsection}{7.07.6 Maintenance, Replacement, and Storage}

\begin{itemize}
\item Inspection cadence:
  \begin{itemize}
  \item pre-use quick check: buckles, straps, stitching, load module closures.
  \item periodic deeper check: frame integrity, hip belt stitching, compression strap anchors.
  \end{itemize}
\item Cleaning/drying:
  \begin{itemize}
  \item sweat management and mildew prevention: dry fully after use.
  \item storage dry and ventilated; avoid compressing frames in a way that deforms structure.
  \end{itemize}
\item Consumables/spares:
  \begin{itemize}
  \item strap keepers, spare buckle(s), tape, sandbag liners, seam repair tape (generic).
  \end{itemize}
\item Storage footprint:
  \begin{itemize}
  \item store packs and load modules to avoid trip hazards.
  \item avoid storing heavy loads on soft flooring that can deform surfaces; use stable surfaces.
  \end{itemize}
\end{itemize}

\subsection*{7.07.7 Substitution Rules}
\phantomsection
\addcontentsline{toc}{subsection}{7.07.7 Substitution Rules}

Tier 0 substitution rule (non-negotiable):

\begin{itemize}
\item Tier 0 means no load carriage capability. No ruck substitutions are permitted. Do not use household backpacks or improvised loads.
\end{itemize}

Tier 1 and higher substitution posture:

\begin{itemize}
\item Substitutions must preserve load stability, fit integrity, and load verification posture.
\item Permitted substitution examples (category-level):
  \begin{itemize}
  \item pack model substitution only after stabilization period outside protected windows,
  \item load module substitution only if it is non-leaking, non-shifting, and can be labeled and logged.
  \end{itemize}
\item Prohibited substitutions:
  \begin{itemize}
  \item unstable loads that shift,
  \item improvised loads likely to leak/break,
  \item excessive loads without fit and retention integrity,
  \item carrying water as “dry load” (water is logged separately).
  \end{itemize}
\end{itemize}

Validity/logging linkage (Stage 2.E posture, high level; required fields conceptually):

\begin{itemize}
\item When any substitution occurs, log:
  \begin{itemize}
  \item environment identity (Route \textbf{ID} / environment unit posture concept),
  \item tool/load verification method (cross-reference 7.18),
  \item dry load recorded; carried water logged separately,
  \item any deviation reason and the applicable validity posture (if the change affects comparability, treat as validity-sensitive and reslot rather than claim gate credit).
  \end{itemize}
\end{itemize}

\subsection*{7.07.8 Phase Integration Map}
\phantomsection
\addcontentsline{toc}{subsection}{7.07.8 Phase Integration Map}

\begingroup
\scriptsize
\setlength{\tabcolsep}{2pt}
\renewcommand{\arraystretch}{1.12}

\begin{longtblr}{
  width=\linewidth,
  rowhead=1,
  colspec = {
    Q[l,wd=1.05in]
    Q[l,wd=1.60in]
    X[l,3.05]
    X[l,3.95]
  },
  hlines,
  vlines,
  cells = {hsep=6pt, vsep=6pt, halign=l, valign=t},
  row{1} = {bg=StageBlue!15, fg=black, font=\bfseries, halign=l, valign=m},
  row{odd} = {bg=white},
  row{even} = {bg=LightGray},
}
\Hdr{Phase} &
\Hdr{Required Tier (from Stage 6)} &
\Hdr{What Changes / Becomes Mandatory} &
\Hdr{Notes} \\
Phase 00 &
Tier 0 &
No ruck capability; no substitutions permitted &
Phase 00 is minimal. No load carriage requirement. \\
Phase 01 &
Tier 0 &
No ruck capability; education-only posture &
Maintain Tier 0. Do not treat “hiking with a backpack” as authorized ruck progression. \\
Phase 02 &
Tier 0 &
No ruck capability; maintain prohibition on substitutions &
Ruck interface remains out of scope for minimum deployable readiness. \\
Phase 03 &
Tier 1 &
First minimum deployable ruck capability becomes required by Phase 03 (End) &
Milestone: first structured ruck exposure becomes feasible only after Tier 1 capability exists. Notes: Phase 03 (End) is the first required point; stabilize the interface in earlier work weeks; no novelty in protected windows. \\
Phase 04 &
Tier 1 &
Tier 1 ruck capability must be stable by Phase 04 (E) for phase execution &
Notes: Phase 04 (E) indicates ruck capability must be available at phase entry; treat fit and load module selection as locked variables for the phase to protect comparability. \\
Phase 05 &
Tier 1 &
Maintain Tier 1; prepare for higher consequence by stabilizing foot and sustainment posture &
Notes: Phase 05 (End) remains Tier 1; do not escalate interface changes near longer endurance gates. \\
Phase 06 &
Tier 2 &
Tier 2 expanded stability becomes required by Phase 06 (End) &
Notes: Phase 06 (End) requires improved stability and comfort posture; prioritize repeatability, not load escalation. \\
Phase 07 &
Tier 2 &
Maintain Tier 2 stability under hybrid density &
Notes: Phase 07 (End) keep interface stable across protected windows; reslot rather than compress if variance occurs. \\
Phase 08 &
Tier 2 &
Maintain Tier 2; load tolerance consequence increases &
Notes: Phase 08 (End) treat long-load milestones as dominant events by intent; avoid interface changes near protected windows. \\
Phase 09 &
Tier 2 &
Maintain Tier 2; manage friction and hotspots under volume &
Notes: Phase 09 (End) \textbf{CAL} volume and ruck exposure must not collide through avoidable foot breakdown; cross-reference 7.11. \\
Phase 10 &
Tier 2 &
Maintain Tier 2; multi-domain gate weeks increase variance sensitivity &
Notes: Phase 10 (End) high consequence; enforce no novelty posture; visibility and environment controls become decisive. \\
Phase 11 &
Tier 2 &
Maintain Tier 2; recovery sensitivity increases &
Notes: Phase 11 (End) maintain stable interface; reslot if risk flags appear; do not force. \\
Phase 12 &
Tier 3 &
Tier 3 redundancy becomes required by Phase 12 (End) per Stage 6 &
Notes: Phase 12 (End) redundancy is introduced as reliability posture for high-consequence windows; stabilize any new pack/frame well before verification blocks. \\
Phase 13 &
Tier 3 &
Maintain Tier 3 stability posture through consolidation &
Notes: Phase 13 (End) no novelty near final protected windows; repeatability is the objective. \\
\end{longtblr}

\endgroup

7.07.8 Notes:

\begin{itemize}
\item If any upstream phase artifact implies ruck capability is required before Phase 03 (End), Requires Stage 8 review.
\item If any upstream phase artifact implies underwater/breath-hold is a substitute for load carriage readiness, Requires Stage 8 review.
\end{itemize}

\subsection*{7.07.9 Risks, Contraindications, and Red Flags}
\phantomsection
\addcontentsline{toc}{subsection}{7.07.9 Risks, Contraindications, and Red Flags}

Non-medical risk list (Phase-gated posture):

\begin{itemize}
\item Foot hotspots/blisters (cross-reference 7.11):
  \begin{itemize}
  \item hotspot trending toward blister requires downshift next exposure and stabilization before escalation.
  \item active blister increases fall risk and invalidates gait comparability; substitute lower-risk Cardio and reslot ruck exposures.
  \end{itemize}
\item Knee/hip/back flare-ups:
  \begin{itemize}
  \item gait-altering pain triggers stop posture for load carriage; substitute lower-risk modality; document.
  \end{itemize}
\item Heat illness risk:
  \begin{itemize}
  \item increased risk under load; enforce conservative route selection and hydration posture (7.13); stop for heat illness signs.
  \end{itemize}
\item Traffic/visibility hazards:
  \begin{itemize}
  \item low light and traffic elevate consequence; visibility posture (7.19) is controlling; reslot if safe separation cannot be maintained.
  \end{itemize}
\item Load shift and fall risk:
  \begin{itemize}
  \item shifting loads increase trip/fall risk; any shift indicates the system is not deployable at that load level; stop and correct.
  \end{itemize}
\item Red-flag symptom posture reminder:
  \begin{itemize}
  \item chest pain/pressure, syncope/near-syncope, disproportionate dyspnea, new neurologic symptoms: \textbf{STOP} \textbf{NOW} and seek evaluation.
  \item gait-altering pain: stop ruck exposure; substitute lower-risk Cardio; document per Stage 2.E posture conceptually.
  \end{itemize}
\item Timing mismatch risk:
  \begin{itemize}
  \item If Stage 6 timing and any phase artifact conflict on when Tier 1 is required, flag Requires Stage 8 review rather than changing timing here.
  \end{itemize}
\end{itemize}

\end{document}